%*************************************
% บทที่ 1
%*************************************
\chapterTitle{1}{บทนำ}

% คำสั่งสำหรับจัดรายการหัวข้อย่อยแบบชิดสวยงาม ไม่มีช่องว่างกว้างผิดปกติ
\providecommand{\decitem}[2]{%
  \par\noindent\hangindent=3.6em\hangafter=1\hspace*{1.3em}#1\hspace{0.5em}#2\par\vspace{0.25em}%
}
\providecommand{\subdecitem}[2]{%
  \par\noindent\hangindent=5.8em\hangafter=1\hspace*{2.0em}#1\hspace{0.5em}#2\par\vspace{0.22em}%
}

\section{ความเป็นมาและความสำคัญของปัญหา}
\hspace*{1.3em}
ในยุคเศรษฐกิจดิจิทัลปัจจุบัน องค์กรธุรกิจขนาดเล็กและขนาดกลาง (Small and Medium Enterprises: SMEs) มีบทบาทสำคัญอย่างยิ่งต่อการขับเคลื่อนเศรษฐกิจของประเทศ อย่างไรก็ดี องค์กรขนาดเล็กถึงขนาดกลางมักเผชิญกับข้อจำกัดด้านทรัพยากรบุคคล งบประมาณ และเครื่องมือทางเทคโนโลยีที่ใช้ในการบริหารจัดการข้อมูล การดำเนินงานส่วนใหญ่ยังคงพึ่งพาการทำงานแบบดั้งเดิม (Manual Process) เช่น การบันทึกบัญชีลงสมุด การจดบันทึกเอกสารลงสเปรดชีตที่กระจัดกระจายตามเครื่องของพนักงานแต่ละคน และการประสานงานผ่านแอปพลิเคชันสนทนาที่ไม่มีการจัดหมวดหมู่อย่างเป็นระบบ ส่งผลให้เกิดปัญหาข้อมูลสูญหาย การสื่อสารคลาดเคลื่อน และความล่าช้าในการดำเนินธุรกิจ

\hspace*{1.3em}
กรณีศึกษา บริษัท ฮาร์ทวอมมิ่ง จำกัด ซึ่งเป็นผู้นำเข้าและตัวแทนจัดจำหน่ายสินค้าไลฟ์สไตล์พรีเมียมจากต่างประเทศอย่างเป็นทางการในประเทศไทย ได้แก่ แบรนด์ "duit" ผลิตภัณฑ์นวัตกรรมของใช้และเฟอร์นิเจอร์สัตว์เลี้ยงดีไซน์มินิมอล เช่น เครื่องให้อาหารอัตโนมัติ น้ำพุสัตว์เลี้ยง และกระบะทรายอัจฉริยะ, แบรนด์ "thumb toe" ถุงเท้าพิลาทิสและถุงเท้าแฟชั่นออกกำลังกายที่เน้นการทรงตัว และแบรนด์ "by.tuneteller" ฟองน้ำแต่งหน้าแบบตบแห้งคุณภาพสูง โดยมีช่องทางการจัดจำหน่ายหลักผ่านแพลตฟอร์มอีคอมเมิร์ซและมาร์เก็ตเพลสชั้นนำ ได้แก่ Shopee, Lazada, TikTok Shop และ LINE Official Account (LINE OA) ควบคู่กับการออกบูธแสดงสินค้าในงานนิทรรศการต่าง ๆ เพื่อสร้างการรับรู้แบรนด์

\hspace*{1.3em}
จากการศึกษาและวิเคราะห์กระบวนการทำงานจริงภายใน บริษัท ฮาร์ทวอมมิ่ง จำกัด พบว่าองค์กรกำลังเผชิญกับปัญหาสำคัญในการดำเนินงาน 4 ประการหลัก ดังนี้:

\vspace{0.2em}
\noindent\hspace*{1.3em}\textbf{1) ปัญหาการค้นหาข้อมูลและเอกสารองค์กร:} ข้อมูลรายละเอียดสินค้า คุณสมบัติ คู่มือการใช้งาน นโยบายการรับประกันสินค้า ระเบียบการปฏิบัติงาน และคำถามที่พบบ่อย (FAQ) ถูกจัดเก็บแยกย่อยในเอกสารไฟล์ PDF, ข้อความ และสเปรดชีตหลายแหล่ง เมื่อพนักงานต้องการสืบค้นข้อมูลเพื่อตอบคำถามลูกค้าหรือแก้ไขปัญหาหน้างาน ต้องใช้เวลาค้นหานานและบ่อยครั้งที่พนักงานใหม่ให้ข้อมูลแก่ลูกค้าไม่ถูกต้องตรงกัน\par\vspace{0.3em}

\noindent\hspace*{1.3em}\textbf{2) ปัญหาความผิดพลาดและความล่าช้าในการทำบัญชีรายรับ-รายจ่าย:} ในแต่ละวันมีรายการรับเงินและจ่ายเงินผ่านช่องทางออนไลน์จำนวนมาก พนักงานต้องรวบรวมภาพถ่ายใบเสร็จรับเงิน ใบกำกับภาษี และสลิปการโอนเงินของลูกค้าและซัพพลายเออร์ จากนั้นต้องพิมพ์กรอกตัวเลขและรายละเอียดลงในสมุดบัญชีด้วยมือ (Manual Data Entry) ซึ่งใช้เวลามาก เกิดความผิดพลาดของตัวเลข และเกิดความล่าช้าในการสรุปงบการเงินประจำเดือน\par\vspace{0.3em}

\noindent\hspace*{1.3em}\textbf{3) ปัญหาการบริหารงานบุคคลที่ขาดระบบศูนย์กลาง:} การยื่นคำขอลางานของพนักงาน เช่น ลาป่วย ลากิจ ลาพักร้อน ยังใช้วิธีการส่งข้อความส่วนตัวหรือกระดาษ ทำให้หัวหน้างานตรวจสอบประวัติการลาและคำนวณโควตาวันลาคงเหลือได้ยาก ข้อมูลการขาด ลา มาสาย มักตกหล่นและเกิดความสับสนในการบันทึกสถิติการทำงาน\par\vspace{0.3em}

\noindent\hspace*{1.3em}\textbf{4) ปัญหาความยุ่งยากจากการสลับใช้งานหลายซอฟต์แวร์ (App Fatigue):} พนักงานต้องสลับการทำงานไปมาระหว่างหลายแอปพลิเคชัน เช่น โปรแกรมแชทสำหรับคุยกับลูกค้า สเปรดชีตสำหรับเช็คสต็อก โฟลเดอร์รวมสลิปการเงิน และแอปพลิเคชันสำหรับยื่นใบลา ซึ่งนอกจากจะสร้างความยุ่งยากในการเรียนรู้ระบบแล้ว การจัดซื้อซอฟต์แวร์สำเร็จรูปหลายตัวยังมีค่าบริการรายเดือนที่สูงเกินไปสำหรับธุรกิจ SME\par\vspace{0.4em}

\hspace*{1.3em}
จากสภาพปัญหาดังกล่าว คณะผู้จัดทำจึงได้เล็งเห็นถึงโอกาสในการนำเทคโนโลยีปัญญาประดิษฐ์ (Artificial Intelligence: AI) และโมเดลภาษาขนาดใหญ่ (Large Language Models: LLM) มาบูรณาการร่วมกับแพลตฟอร์ม LINE ซึ่งเป็นแอปพลิเคชันสนทนาที่พนักงานในประเทศไทยมีความคุ้นเคยและใช้งานเป็นประจำอยู่แล้ว ผ่านเทคโนโลยี LINE Messaging API ร่วมกับการพัฒนาเว็บแอปพลิเคชันบริหารจัดการส่วนกลาง (Centralized Web Dashboard) โดยประยุกต์ใช้เทคนิคการค้นคืนข้อมูลเสริมด้วยปัญญาประดิษฐ์ (Retrieval-Augmented Generation: RAG) เพื่อสร้างระบบคลังความรู้อัจฉริยะที่สามารถสืบค้นคำตอบจากเอกสารของบริษัทได้อย่างแม่นยำ ผสานกับเทคโนโลยีการรู้จำอักขระด้วยแสง (Optical Character Recognition: OCR) ในการอ่านและสกัดข้อมูลจากภาพถ่ายสลิปโอนเงินและใบเสร็จ แล้วนำไปลงบัญชีใน Google Sheets พร้อมจัดเก็บเอกสารลง Google Drive ให้อัตโนมัติ ตลอดจนพัฒนาระบบบริหารงานบุคคลสำหรับการขอลางานผ่าน LINE ได้อย่างสะดวกรวดเร็ว

\hspace*{1.3em}
การพัฒนาระบบปัญญาประดิษฐ์แชทบอทสำหรับองค์กรนี้ จึงไม่เพียงแต่ช่วยลดขั้นตอนและระยะเวลาในการทำงานของพนักงาน แต่ยังช่วยลดข้อผิดพลาดในการลงบัญชี เพิ่มความคล่องตัวในการดำเนินธุรกิจ และช่วยให้ผู้บริหารมีข้อมูลสารสนเทศที่ถูกต้อง ครบถ้วน เพื่อนำไปใช้ประกอบการตัดสินใจทางธุรกิจได้อย่างมีประสิทธิภาพและคุ้มค่าต่อการลงทุน

\section{วัตถุประสงค์ของการจัดทำโครงงานพิเศษ}
\hspace*{1.3em}
การจัดทำโครงงานพิเศษนี้ มีวัตถุประสงค์หลัก 5 ประการ ดังนี้:

\vspace{0.2em}
\decitem{1.2.1}{เพื่อวิเคราะห์ ออกแบบ และพัฒนาระบบแชทบอทอัจฉริยะที่ทำงานผ่านแพลตฟอร์ม LINE Official Account ร่วมกับเว็บแอปพลิเคชัน สำหรับช่วยเหลืองานและบริหารจัดการภายในองค์กร}

\decitem{1.2.2}{เพื่อประยุกต์ใช้เทคโนโลยีปัญญาประดิษฐ์ในการสร้างระบบคลังความรู้อัจฉริยะ (Retrieval-Augmented Generation: RAG) ที่สามารถค้นหาคำตอบและดึงข้อมูลจากเอกสารองค์กรได้อย่างถูกต้องและอัตโนมัติ}

\decitem{1.2.3}{เพื่อพัฒนาระบบจัดการการเงินและบัญชีที่ใช้เทคโนโลยี OCR ในการอ่านข้อมูลจากใบเสร็จรับเงินและสลิปโอนเงิน แล้วนำข้อมูลไปบันทึกลงบัญชีใน Google Sheets พร้อมจัดเก็บไฟล์บน Google Drive โดยอัตโนมัติ}

\decitem{1.2.4}{เพื่อพัฒนาระบบบริหารงานบุคคลเบื้องต้น เช่น ระบบการยื่นคำขอลางาน ระบบการอนุมัติใบลาผ่านช่องทาง LINE และระบบตรวจสอบสิทธิ์วันลาคงเหลือของพนักงาน}

\decitem{1.2.5}{เพื่อพัฒนาเว็บแอปพลิเคชันสำหรับผู้ดูแลระบบ (Admin Web Dashboard) ที่แสดงผลสรุปภาพรวมการดำเนินงานขององค์กร การจัดการผู้ใช้งาน การควบคุมสิทธิ์ตามบทบาท (RBAC) และบันทึกประวัติการทำงานของระบบได้อย่างครบถ้วน}

\section{ขอบเขตของการทำโครงงานพิเศษ} \label{sec1:scope}
\hspace*{1.3em}
การดำเนินงานจัดทำโครงงานพิเศษ แบ่งการดำเนินงานออกเป็น 2 ภาคการศึกษา โดยมีขอบเขตของระบบครอบคลุมรายละเอียดดังนี้:

\subsection{ภาคการศึกษาที่ 1}
\hspace*{1.3em}
โครงสร้างพื้นฐานระบบจัดการผู้ใช้ ระบบสมาชิก คลังความรู้ RAG และการจัดการไฟล์

\vspace{0.4em}
\noindent\hspace*{1.3em}\textbf{1.3.1.1 ด้านระบบจัดการผู้ใช้งาน (User Management)}\par\nopagebreak\vspace{0.25em}
\subdecitem{1.3.1.1.1}{สมัครสมาชิกและยืนยันตัวตนเข้าสู่ระบบผ่านอีเมลและรหัสผ่าน (Email \& Password Authentication)}
\subdecitem{1.3.1.1.2}{เข้าสู่ระบบผ่านบัญชี LINE (LINE Login OAuth 2.0) และบัญชี Google (Google OAuth 2.0)}
\subdecitem{1.3.1.1.3}{มีระบบจัดการข้อมูลโปรไฟล์ส่วนตัวของผู้ใช้งาน การแก้ไขรหัสผ่าน และการกู้คืนรหัสผ่านด้วย OTP}
\subdecitem{1.3.1.1.4}{มีระบบควบคุมความปลอดภัยของเซสชัน (Session Management) และตรวจสอบสถานะการเข้าสู่ระบบ}

\vspace{0.4em}
\noindent\hspace*{1.3em}\textbf{1.3.1.2 ด้านระบบแชทบอทอัจฉริยะและคลังความรู้ (AI Chatbot \& Knowledge Base)}\par\nopagebreak\vspace{0.25em}
\subdecitem{1.3.1.2.1}{เชื่อมต่อระบบแชทบอทกับบัญชี LINE Official Account ขององค์กรผ่าน LINE Messaging API}
\subdecitem{1.3.1.2.2}{ระบบคลังความรู้อัจฉริยะ (RAG System) ที่รองรับการอัปโหลดไฟล์เอกสารองค์กร ได้แก่ PDF, DOCX, TXT และสเปรดชีต}
\subdecitem{1.3.1.2.3}{ระบบแปลงเอกสารเป็นเวกเตอร์เอ็มเบดดิง (Vector Embeddings) และจัดเก็บในฐานข้อมูลเวกเตอร์ ChromaDB}
\subdecitem{1.3.1.2.4}{ระบบสืบค้นข้อมูลตามบริบทความหมาย (Semantic Similarity Search) เพื่อดึงข้อความที่เกี่ยวข้องมาสร้างคำตอบด้วยโมเดลภาษาขนาดใหญ่ Google Gemini}

\vspace{0.4em}
\noindent\hspace*{1.3em}\textbf{1.3.1.3 ด้านระบบจัดการเอกสารองค์กร (Document Management)}\par\nopagebreak\vspace{0.25em}
\subdecitem{1.3.1.3.1}{รองรับการอัปโหลดไฟล์หลายรูปแบบผ่านหน้าเว็บและช่องแชท LINE}
\subdecitem{1.3.1.3.2}{เชื่อมต่อกับบริการคลาวด์ Google Drive API เพื่อจัดระเบียบไฟล์แยกโฟลเดอร์ตามวันที่และหมวดหมู่อัตโนมัติ}
\subdecitem{1.3.1.3.3}{รองรับการสืบค้นและดึงลิงก์ไฟล์ย้อนหลังผ่านคำสั่งในช่องแชท LINE}

\vspace{0.4em}
\noindent\hspace*{1.3em}\textbf{1.3.1.4 ด้านระบบเชื่อมต่อแพลตฟอร์ม LINE (LINE Platform Integration)}\par\nopagebreak\vspace{0.25em}
\subdecitem{1.3.1.4.1}{พัฒนาเว็บฮุก (Webhook Handler) สำหรับรับส่งข้อความ เหตุการณ์ และมัลติมีเดีย (รูปภาพ, เอกสาร, พิกัดตำแหน่ง) อัตโนมัติ}
\subdecitem{1.3.1.4.2}{ออกแบบข้อความตอบกลับในรูปแบบการ์ดข้อมูลยืดหยุ่น (LINE Flex Message) ที่สวยงาม อ่านง่าย และมีปุ่มปฏิสัมพันธ์}

\vspace{0.5em}
\subsection{ภาคการศึกษาที่ 2}
\hspace*{1.3em}
ระบบการเงิน OCR, บริหารงานบุคคล, แผงควบคุม และการทดสอบ

\vspace{0.4em}
\noindent\hspace*{1.3em}\textbf{1.3.2.1 ด้านระบบจัดการการเงินและบัญชีอัตโนมัติ (Financial \& Bookkeeping Automation)}\par\nopagebreak\vspace{0.25em}
\subdecitem{1.3.2.1.1}{รองรับการรับภาพถ่ายสลิปโอนเงินและใบเสร็จรับเงินผ่านช่องแชท LINE}
\subdecitem{1.3.2.1.2}{ประยุกต์ใช้เทคโนโลยี OCR (Google Cloud Vision API / Tesseract) ร่วมกับเทคนิค Regular Expression ในการสกัดข้อมูลสำคัญ เช่น วันที่ เวลา จำนวนเงิน ผู้รับโอน/ผู้โอน ธนาคาร และรหัสอ้างอิง}
\subdecitem{1.3.2.1.3}{ระบบนำข้อมูลทางการเงินที่สกัดได้ไปบันทึกลงใน Google Sheets โดยแยกหมวดหมู่รายรับ-รายจ่ายให้อัตโนมัติ}
\subdecitem{1.3.2.1.4}{ระบบกระทบยอดอัจฉริยะ (Reconciliation Engine) สำหรับจับคู่สลิปโอนเงินกับเอกสารใบกำกับภาษีหรือใบสั่งซื้อ}

\vspace{0.4em}
\noindent\hspace*{1.3em}\textbf{1.3.2.2 ด้านระบบบริหารงานบุคคล (Human Resource Management)}\par\nopagebreak\vspace{0.25em}
\subdecitem{1.3.2.2.1}{ระบบยื่นใบลาผ่านช่องแชท LINE รองรับการลาป่วย ลากิจ ลาพักร้อน พร้อมแนบเอกสารรับรอง}
\subdecitem{1.3.2.2.2}{ระบบส่งการแจ้งเตือนแบบเรียลไทม์ไปยังหัวหน้างานทาง LINE เพื่อพิจารณาอนุมัติหรือปฏิเสธคำขอได้ทันที}
\subdecitem{1.3.2.2.3}{ระบบตรวจสอบและคำนวณโควตาวันลาคงเหลือของพนักงานแต่ละคนโดยอัตโนมัติ}
\subdecitem{1.3.2.2.4}{ระบบลงเวลาเข้า-ออกงาน (Time Attendance) พร้อมบันทึกพิกัดตำแหน่ง (Geolocation) ผ่าน LINE}

\newpage
\noindent\hspace*{1.3em}\textbf{1.3.2.3 ด้านระบบแผงควบคุมสำหรับผู้บริหารและผู้ดูแลระบบ (Executive \& Admin Dashboard)}\par\nopagebreak\vspace{0.25em}
\subdecitem{1.3.2.3.1}{แผงควบคุมแสดงผลสถิติและกราฟภาพรวม (Data Visualization) ยอดรายรับ-รายจ่าย สถิติการใช้งาน AI และข้อมูลการลา}
\subdecitem{1.3.2.3.2}{ระบบจัดการสมาชิก กำหนดระดับสิทธิ์การเข้าถึงข้อมูลตามบทบาท (Role-Based Access Control: RBAC)}
\subdecitem{1.3.2.3.3}{ระบบจัดการกฎการจำแนกหมวดหมู่บัญชีของ AI (AI Category Rules)}
\subdecitem{1.3.2.3.4}{ระบบตรวจสอบบันทึกการปฏิบัติงาน (Audit Logs) และสถานะการเชื่อมต่อบริการภายนอก (System Health)}
\subdecitem{1.3.2.3.5}{ระบบจัดการเอกสารใบเสนอราคา (Quotations) และระบบจัดการข้อมูลลูกค้าสัมพันธ์ (CRM)}

\vspace{0.4em}
\noindent\hspace*{1.3em}\textbf{1.3.2.4 ด้านระบบรักษาความปลอดภัยและการทดสอบ (Security \& System Evaluation)}\par\nopagebreak\vspace{0.25em}
\subdecitem{1.3.2.4.1}{ระบบแยกข้อมูลระหว่างองค์กร (Multi-Tenant Data Isolation) เพื่อรักษาความปลอดภัยของข้อมูลในระดับองค์กร}
\subdecitem{1.3.2.4.2}{ระบบตรวจสอบความถูกต้องของไฟล์ด้วยการคำนวณค่าแฮช SHA-256 (File Integrity Record)}
\subdecitem{1.3.2.4.3}{ดำเนินการทดสอบระบบเชิงหน้าที่ (Functional Testing), ทดสอบประสิทธิภาพความแม่นยำ OCR และ RAG ตลอดจนประเมินความพึงพอใจของผู้ใช้งานจริง}

\vspace{0.5em}
\section{ประโยชน์ที่คาดว่าจะได้รับ}
\vspace{0.2em}
\decitem{1.4.1}{ช่วยลดขั้นตอน ระยะเวลา และภาระงานของพนักงาน บริษัท ฮาร์ทวอมมิ่ง จำกัด ในการค้นหาข้อมูลสินค้า การตอบคำถามลูกค้า และการค้นหาเอกสารองค์กร}
\decitem{1.4.2}{ช่วยเพิ่มความถูกต้องและลดความผิดพลาดในการให้ข้อมูลสินค้าแก่ลูกค้า โดยมีระบบ AI RAG คอยอ้างอิงจากคลังเอกสารและคู่มือมาตรฐานของบริษัท}
\decitem{1.4.3}{ช่วยให้การจัดเก็บเอกสารและไฟล์ภาพขององค์กรมีความเป็นระเบียบ เรียบร้อย และปลอดภัยผ่านการเชื่อมต่อบริการคลาวด์ Google Drive}
\decitem{1.4.4}{ช่วยประหยัดค่าใช้จ่ายขององค์กร เนื่องจากเป็นการพัฒนาขึ้นมาใช้งานเองภายในองค์กรโดยอาศัยแพลตฟอร์ม LINE ที่พนักงานคุ้นเคย โดยไม่ต้องเสียค่าบริการซอฟต์แวร์สำเร็จรูปรายเดือนที่มีราคาสูง}
\decitem{1.4.5}{สามารถนำองค์ความรู้และสถาปัตยกรรมระบบไปประยุกต์ใช้เป็นต้นแบบในการพัฒนาระบบผู้ช่วยอัจฉริยะสำหรับธุรกิจขนาดกลางและขนาดย่อม (SMEs) อื่น ๆ ได้ในอนาคต}

\section{ผลที่คาดว่าจะได้รับ}
\vspace{0.2em}
\decitem{1.5.1}{ได้ระบบผู้ช่วยอัจฉริยะประจำองค์กรที่ทำงานผ่าน LINE Official Account ร่วมกับเว็บแอปพลิเคชัน ซึ่งช่วยให้พนักงานสามารถสืบค้นข้อมูล ค้นหาเอกสาร ยื่นคำขอลางาน และบันทึกข้อมูลการเงินได้จากแพลตฟอร์มเดียว โดยไม่ต้องติดตั้งแอปพลิเคชันเพิ่มเติม}
\decitem{1.5.2}{ได้ระบบคลังความรู้อัจฉริยะ (RAG) ที่สามารถตอบคำถามเกี่ยวกับข้อมูลสินค้า นโยบาย และระเบียบปฏิบัติของ บริษัท ฮาร์ทวอมมิ่ง จำกัด ได้อย่างแม่นยำและถูกต้อง ช่วยลดเวลาในการค้นหาข้อมูลของพนักงานลงได้มากกว่าร้อยละ 70}
\decitem{1.5.3}{ได้ระบบประมวลผลเอกสารทางการเงินอัตโนมัติด้วย OCR ที่สามารถอ่านข้อมูลจากภาพถ่ายใบเสร็จและสลิปโอนเงิน แล้วลงบัญชีใน Google Sheets พร้อมจัดเก็บเอกสารบน Google Drive ช่วยลดข้อผิดพลาดในการคีย์ข้อมูลด้วยมือ และลดเวลาการทำบัญชีประจำวัน}
\decitem{1.5.4}{ได้ระบบบริหารงานบุคคลที่สะดวกรวดเร็ว โดยพนักงานสามารถยื่นใบลาและตรวจสอบโควตาวันลาผ่าน LINE ได้ทันที ขณะที่หัวหน้างานสามารถพิจารณาอนุมัติผ่าน LINE Flex Message ได้แบบเรียลไทม์}
\decitem{1.5.5}{ได้แผงควบคุมผู้บริหาร (Admin Dashboard) ที่สรุปสถิติภาพรวมการดำเนินงานขององค์กร รายรับ-รายจ่าย การลา และรายงานสารสนเทศ ซึ่งช่วยให้ผู้บริหารสามารถติดตามสถานะการดำเนินงานและตัดสินใจเชิงธุรกิจได้อย่างทันท่วงที}

\section{ข้อตกลงเบื้องต้น}
\vspace{0.2em}
\decitem{1.6.1}{ข้อมูลเอกสาร คู่มือสินค้า และระเบียบปฏิบัติที่นำมาใช้ทดสอบคลังความรู้ เป็นข้อมูลจริงของ บริษัท ฮาร์ทวอมมิ่ง จำกัด ได้แก่ แบรนด์ duit, thumb toe และ by.tuneteller}
\decitem{1.6.2}{ระบบทำงานผ่านการเชื่อมต่อเครือข่ายอินเทอร์เน็ตความเร็วสูง และพึ่งพาการทำงานของภายนอก เช่น LINE Platform, Google Cloud Services (Gemini API, Google Drive API)}
\decitem{1.6.3}{ความถูกต้องในการสืบค้นข้อมูลของระบบ RAG ขึ้นอยู่กับคุณภาพ ความสมบูรณ์ และการจัดหมวดหมู่ของเอกสารต้นฉบับที่นำเข้าสู่ระบบ}
\decitem{1.6.4}{การดำเนินงานในภาคการศึกษาที่ 1/2569 มุ่งเน้นการวิเคราะห์ ออกแบบระบบ และพัฒนาระบบโครงสร้างพื้นฐาน ได้แก่ ระบบสมาชิก, ระบบเชื่อมต่อ LINE OA, ระบบคลังความรู้ RAG, และระบบจัดการไฟล์ Google Drive สำหรับระบบการเงิน OCR, ระบบบริหารงานบุคคล และการทดสอบระบบเต็มรูปแบบจะดำเนินการในภาคการศึกษาที่ 2/2569}

\section{นิยามศัพท์เฉพาะ}
\vspace{0.2em}
\decitem{1.7.1}{Retrieval-Augmented Generation (RAG) หมายถึง สถาปัตยกรรมทางปัญญาประดิษฐ์ที่ผสมผสานการค้นคืนสารสนเทศจากคลังเอกสารเฉพาะขององค์กรเข้ากับความสามารถในการสร้างภาษาธรรมชาติของโมเดลภาษาขนาดใหญ่ เพื่อให้คำตอบมีความถูกต้องตรงตามข้อเท็จจริง}
\decitem{1.7.2}{Large Language Model (LLM) หมายถึง แบบจำลองภาษาขนาดใหญ่ที่ถูกฝึกฝนด้วยข้อมูลมหาศาล มีความสามารถในการเข้าใจความหมายของภาษาธรรมชาติ ประมวลผลข้อความ และสร้างคำตอบได้อย่างสละสลวยและเป็นธรรมชาติ}
\decitem{1.7.3}{Vector Embeddings หมายถึง การแปลงข้อมูลข้อความหรือเอกสารให้อยู่ในรูปเวกเตอร์ตัวเลขหลายมิติ เพื่อให้คอมพิวเตอร์สามารถวัดระยะห่างและเปรียบเทียบความหมายเชิงบริบท (Semantic Similarity) ได้อย่างแม่นยำ}
\decitem{1.7.4}{ChromaDB หมายถึง ระบบจัดการฐานข้อมูลเวกเตอร์แบบโอเพนซอร์ส (Vector Database) ที่ออกแบบมาเพื่อจัดเก็บและค้นคืนเวกเตอร์เอ็มเบดดิงที่มีประสิทธิภาพสูง}
\decitem{1.7.5}{LINE Messaging API หมายถึง ชุดบริการเชื่อมต่อระหว่างเซิร์ฟเวอร์ของระบบกับแพลตฟอร์ม LINE ช่วยให้สามารถรับส่งข้อความ เหตุการณ์ และมัลติมีเดียแบบสองทางอัตโนมัติ}
\decitem{1.7.6}{LINE Flex Message หมายถึง รูปแบบข้อความตอบกลับบน LINE ที่สามารถจัดวางองค์ประกอบ ภาพ ตัวอักษร และปุ่มกดได้อย่างยืดหยุ่นและสวยงามตามหลักการออกแบบส่วนต่อประสาน}
\decitem{1.7.7}{Webhook หมายถึง กลไกการสื่อสารระหว่างระบบผ่าน HTTP POST request ที่ส่งข้อมูลเหตุการณ์ไปยังเซิร์ฟเวอร์ปลายทางทันทีที่เกิดเหตุการณ์ขึ้นแบบเรียลไทม์}
\decitem{1.7.8}{Role-Based Access Control (RBAC) หมายถึง ระบบการควบคุมการเข้าถึงข้อมูลและฟังก์ชันการทำงานของระบบตามบทบาทหน้าที่ของผู้ใช้งาน เช่น ผู้ดูแลระบบ (Admin) หรือพนักงานทั่วไป (Employee)}
\decitem{1.7.9}{Multi-Tenant Data Isolation หมายถึง หลักการออกแบบสถาปัตยกรรมระบบและฐานข้อมูลที่แบ่งแยกข้อมูลระหว่างแต่ละองค์กรหรือหน่วยงานอย่างเด็ดขาด เพื่อความปลอดภัยและความเป็นส่วนตัวของสารสนเทศ}
\decitem{1.7.10}{Optical Character Recognition (OCR) หมายถึง เทคโนโลยีการรู้จำอักขระด้วยแสงที่สกัดข้อความและตัวเลขจากภาพถ่ายเอกสารดิจิทัลให้กลายเป็นข้อมูลข้อความที่นำไปประมวลผลต่อได้}
